% ============================================================================
% TEST STRATEGY
% ============================================================================

\section{Test Strategy}
\label{section:test-strategy}

This section enumerates every layer of automated testing in the
project, what it verifies, where it runs, and how to invoke it
locally.

\subsection{Layer 1: Firmware host unit tests (Unity)}

\textbf{Where:} \texttt{star-rx72n-firmware/tests/}.
\textbf{Framework:} Unity (single-source vendored).
\textbf{Compiler:} system Clang (\emph{not} GNURX).
\textbf{Run:} \texttt{cmake -B build -S star-rx72n-firmware/tests \&\&
cmake --build build -j 8 \&\& ctest --test-dir build -j 8}, or the
shortcut \texttt{make test-rx72n}.

Sixty test executables, each linking the production source for the
library under test against mock substitutes. Mocks live under
\texttt{tests/mocks/}.

\textbf{Coverage gate:} \texttt{firmware-unit-tests.yml} enforces
100\,\% line / function / branch coverage on the
\texttt{libs/}-filtered subset. Hardware-only sources that cannot
run on the host (\texttt{rx\_usb/*}, \texttt{rx\_crc\_hw.c},
\texttt{rx\_hcsr04} hw/isr/icu, several \texttt{rx\_hal/src} files
that drive registers directly) are explicitly excluded so the gate
honestly reflects what is host-testable.

The \texttt{test\_coverage\_compile\_all} target is a no-op test that
links every host-compilable production file -- its job is to ensure
\texttt{lcov --initial} sees every translation unit so the
denominator of coverage is correct.

\subsection{Layer 2: Cross-build verification}

\textbf{Where:} \texttt{star-rx72n-firmware/CMakeLists.txt}.
\textbf{Compiler:} GNURX 14.2.0 (cross).
\textbf{Run:} \texttt{make build-rx72n} or \texttt{make ci-rx72n}.

This proves the firmware tree compiles for the actual target with
\texttt{-Wall -Wextra -Werror} and links cleanly. It does not run
the firmware -- there is no on-target test harness in this layer.

\textbf{Optional clang-tidy:} \texttt{cmake -DENABLE\_CLANG\_TIDY=ON}
runs the project's curated clang-tidy ruleset. NOLINT suppressions
for \texttt{readability-*} checks are forbidden; failing checks
must be refactored.

\subsection{Layer 3: Gateway Go tests (race + coverage)}

\textbf{Where:} \texttt{star-gateway/internal/}.
\textbf{Framework:} Go's built-in \texttt{testing} + testify.
\textbf{Run:} \texttt{cd star-gateway \&\& go test -race -cover ./...}.

Coverage requirement: $\ge$ 75\,\%. Most packages well above; some
are at 100\,\% (\texttt{controller}, \texttt{fec}, \texttt{frame}).

\textbf{Fuzz tests:} \texttt{frame/fuzz\_test.go},
\texttt{dispatcher/fuzz\_test.go} run under \texttt{go test -fuzz=Fuzz}
on demand and are part of the periodic HIL run.

\subsection{Layer 4: Hardware-in-the-Loop (HIL)}

\textbf{Where:} \texttt{star-gateway/test/e2e/hil\_test.go}.
\textbf{Setup:} Pi 5 + RX72N over USB CDC + AD2 capture rig.
\textbf{Run:} \texttt{go test -tags=hardware ./test/e2e/...}.
\textbf{CI:} \texttt{firmware-hil-nightly.yml} on the self-hosted
Pi 5 runner.

Tests command / response round-trip latency, transport failover,
USB hot-plug, the simple-USB vs full-HARQ mode selection, and
end-to-end estop propagation.

\subsection{Layer 5: ROS2 colcon tests}

\textbf{Where:} \texttt{star-ros2/src/<package>/test/}.
\textbf{Framework:} gtest for C++, pytest for Python.
\textbf{Run:} \texttt{./test-ros2.sh}.

Per-package focus:

\begin{itemize}
  \item \texttt{star\_safety\_monitor}: lifecycle transitions, sonar
        threshold breach, bond failure handling.
  \item \texttt{star\_supervisor}: estop / autonomy gating
        truth-table coverage.
  \item \texttt{star\_spi\_bridge}: frame round-trip with mock
        spidev0.0.
  \item \texttt{star\_gateway\_bridge}: gRPC stub server integration
        tests.
  \item \texttt{star\_compliance\_msgs}: serialization round-trip.
\end{itemize}

\subsection{Layer 6: Compliance engine tests}

\textbf{Where:} \texttt{star-compliance/tests/}.
\textbf{Framework:} pytest.
\textbf{Run:} \texttt{cd star-compliance \&\& PYTHONPATH=. pytest tests/}
or via \texttt{compliance-engine-ci.yml}.

Detector unit tests (cane\_zone\_filter, corridor\_medial\_axis,
doorway\_lidar\_detector, jamb\_plane\_fitter, imu\_jolt\_detector,
obstacle\_clusterer), engine tests (door\_offset\_calibration,
plane\_segmentation, floor\_frame), and node smoke / behavior tests.

\subsection{Layer 7: Proto round-trip tests}

\textbf{Where:} \texttt{star-proto/tests/\{go,typescript,nanopb\}/}.
\textbf{Run:} per-language test commands; CI in \texttt{proto.yml}.

Verifies that a message encoded by one language decodes byte-
identically in every other. \texttt{proto-gen.yml} additionally
fails if regenerating produces a non-empty diff (lockstep gate).

\subsection{Layer 8: Static analysis}

\begin{center}
\begin{tabular}{|p{4.5cm}|p{10.0cm}|}
\hline
\textbf{Tool}                  & \textbf{Scope}                                                                                                                       \\
\hline
clang-format                   & Repo-wide on \texttt{.c/.h/.cpp/.hpp} except ROS2 packages.                                                                            \\
ament\_uncrustify              & ROS2 packages only (\texttt{docs/ROS2\_FORMATTING.md}).                                                                              \\
clang-tidy                     & Optional pass on firmware host build. NOLINT for \texttt{readability-*} is forbidden.                                                  \\
golangci-lint                  & Run from \texttt{gateway.yml}; full suite per \texttt{.golangci.yml}.                                                                \\
govulncheck                    & Go-side vulnerability scanner.                                                                                                        \\
ament\_cppcheck, ament\_cpplint & ROS2 default analyzers.                                                                                                              \\
buf lint                       & Proto schemas; enforced by \texttt{proto.yml}.                                                                                        \\
buf breaking                   & Detects non-additive proto changes.                                                                                                   \\
ASCII purity check             & \texttt{scripts/utils/fix-encoding.py --check .}; CI \texttt{ascii-check.yml}.                                                       \\
Copyright header check         & \texttt{scripts/utils/check-copyright.py}; CI \texttt{copyright-check.yml}.                                                          \\
\texttt{@since} version check  & CI \texttt{since-version-check.yml}.                                                                                                  \\
\hline
\end{tabular}
\end{center}

\subsection{What you should run before every commit}

\begin{enumerate}
  \item Whatever subset of the above touches the code you changed.
  \item For C firmware edits: \texttt{make test-rx72n} at minimum.
  \item For proto edits: \texttt{make proto-gen} and verify the diff
        includes the regenerated outputs.
  \item For ROS2 edits: \texttt{./build-ros2.sh \&\& ./test-ros2.sh}.
  \item For Go edits: \texttt{cd star-gateway \&\& go test -race ./...
        \&\& golangci-lint run}.
  \item Always: \texttt{python3 scripts/utils/fix-encoding.py --check .}.
\end{enumerate}

The repo includes a tracked pre-commit hook at
\texttt{scripts/git/pre-commit} that enforces ASCII purity. Install it
with \texttt{git config core.hooksPath scripts/git} (opt-in;
\texttt{.git/hooks/} is not auto-populated).

\subsection{Coverage gaps to be honest about}

\begin{itemize}
  \item There is no firmware on-target test harness. Cross-build is
        verified, but the firmware itself runs only on hardware and
        the only on-target tests are the bench-test sandboxes
        (\texttt{star-rx72n-firmware/<peripheral>\_test/}) and the
        HIL e2e suite.
  \item There is no soak-test infrastructure that runs for hours.
        The HIL run is a one-shot pass.
  \item There is no automated PCB / electrical test.
\end{itemize}
